\documentclass[aps,onecolumn,nofootinbib,pra]{article}

\usepackage{../../spec_files/arxiv}
\input{../../spec_files/standard_preamble.tex}

\usepackage{import}
\subimport{../../macros/}{all-macros}

\usepackage{minted}

\newcommand{\gausscore}{g}
\newcommand{\gausscoreof}[1]{\gausscore_{#1}}
\newcommand{\gausscoreat}[1]{\gausscore\left[#1\right]}
\newcommand{\gausscoreofat}[2]{\gausscoreof{#1}\left[#2\right]}

\begin{document}
    \title{Tensor network representation of fourier transforms}

    \maketitle
    \date{\today}

    We derive two standard algorithms in the contraction formalism:
    \begin{itemize}
        \item Fast Fourier Transform
        \item Quantum Fourier Transform
    \end{itemize}


    \section{Definition}

    The Fourier transform of vectors in $\mathbb{C}^{\catdim}$ is represented by a matrix
    \begin{align*}
        F[\indexedcatvariable,\headvariable=\headindex]
        = \frac{1}{\sqrt{\catdim}} \expof{2\pi i \cdot \frac{\catindex\cdot\headindex}{\catdim}}
    \end{align*}
    where $\catvariable,\headvariable$ are variables of dimension $\catdim$.
    The Fourier transform of a vector $\vectorat{\catvariable}$ is then the contraction
    \begin{align*}
        \contractionof{\vectorat{\catvariable},F[\catvariable,\headvariable]}{\headvariable} \, .
    \end{align*}

    \subsection{Unfolding}

    We assume $\catdim=2^{\catorder}$ and represent each index $\catindex\in[2^{\catorder}]$ by a tuple $(\catindexof{0},\ldots,\catindexof{\catorder-1})\in[2]^{\catorder}$ such that
    \begin{align*}
        \catindex
        = \sum_{\catenumeratorin} \catindexof{\catenumerator} \cdot 2^{\catorder-\catenumerator-1} \, .
    \end{align*}
    Note that in this convention, $\catindexof{\catorder-1}$ is the least significant bit.
    The unfolded vector $\vectorat{\catvariable}$ is a tensor $\vectorat{\shortcatvariables}$ with indices
    \begin{align*}
        \vectorat{\indexedshortcatvariables}
        = \vectorat{\catvariable=\left(\sum_{\catenumeratorin} \catindexof{\catenumerator} \cdot 2^{\catorder-\catenumerator-1}\right)} \, .
    \end{align*}



    In the unfolded representation the fourier transform is described as the contraction with a tensor $F^{\catorder}[\catvariableof{[\catorder]},\headvariableof{[\catorder]}]$ with coordinates
    \begin{align*}
        F^{\catorder}[\indexedcatvariableof{[\catorder]},\indexedheadvariableof{[\catorder]}]
        = \frac{1}{\sqrt{2^\catorder}} \expof{2\pi i
        \cdot \left(\sum_{\catenumeratorin} \catindexof{\catenumerator} \cdot 2^{\catorder-\catenumerator-1}\right)
        \cdot \left(\sum_{\catenumeratorin} \headindexof{\catenumerator} \cdot 2^{\catorder-\catenumerator-1}\right)
        } \, ,
    \end{align*}
    that is the Fourier transform of $\vectorat{\shortcatvariables}$ is
    \begin{align*}
        \contractionof{\vectorat{\shortcatvariables},F[\shortcatvariables,\headvariableof{[\catorder]}]}{\headvariableof{[\catorder]}} \, .
    \end{align*}


    \section{Tensor network decomposition}

    We now derive a tensor network decomposition of the tensor $F^{\catorder}[\catvariableof{[\catorder]},\headvariableof{[\catorder]}]$.


    To this end we introduce the diagonal unitary
    \begin{align*}
        U^{\catorder}[\indexedshortcatvariables,\headvariableof{[\catorder]}=\headindexof{[\catorder]}]
        = \begin{cases}
              \expof{2\pi i \cdot \left(\sum_{\catenumeratorin} \catindexof{\catenumerator} \cdot 2^{-\catenumerator-2} \right)} & \ifspace \shortcatindices=\headindexof{[\catorder]} \\
              0 & \elsetext
        \end{cases}
    \end{align*}
    Note that $U^{\catorder}$ is diagonal and the tensor product of matrices with variables $\catvariableof{\catenumerator},\headvariableof{\catenumerator}$.
    Diagonality can be used to contract $U^{\catorder}$ to an arbitrary tensor with computational demand in $\mathcal{0}(2^{\catorder})$.
%
%    \stantikzpic{
%
%        \draw (3,-2) -- (3,-1);
%        %\drawvariabledot{3}{-2}
%
%        \drawhwire{0}{2.5}{$\headvariableof{\catorder-1}$}
%        \drawtextnode{1}{1.5}{\colorlabelsize $\vdots$}
%        \drawhwire{0}{-0.5}{$\headvariableof{1}$}
%        \draw (2,-1) rectangle (4,3);
%        \drawtextnode{3}{1}{\corelabelsize $U^{\catorder-1}$}
%        \draw (4,2.5) -- (7,2.5);
%        \drawtextnode{5}{1.5}{\colorlabelsize $\vdots$}
%        \drawhwire{7}{2.5}{$\headvariableof{\catorder-1}$}
%
%        \drawtextnode{10}{0}{$=$}
%
%
%    }


    With these, the Fourier transform on $[2^{\catorder}]$ can be decomposed into a Fourier transform on $[2^{(\catorder-1)}]$ as:
    \stantikzpic{
        \drawhwire{-3}{2.5}{$\catvariableof{0}$}
        \drawtextnode{-2}{0.5}{\colorlabelsize $\vdots$}
        \drawhwire{-3}{-2.5}{$\catvariableof{\catorder-1}$}
        \draw (-1,-3) rectangle (1,3);
        \drawtextnode{0}{0}{\corelabelsize $F^{\catorder}$}
        \drawhwire{1}{2.5}{$\headvariableof{\catorder-1}$}
        \drawtextnode{2}{0.5}{\colorlabelsize $\vdots$}
        \drawhwire{1}{-2.5}{$\headvariableof{0}$}

        \drawtextnode{8}{0}{$=$}

        \begin{scope}[shift={(15,0)}]
            \drawhwire{-4}{2.5}{$\catvariableof{0}$}
            \drawtextnode{-3}{1.5}{\colorlabelsize $\vdots$}
            \drawhwire{-4}{-0.5}{$\catvariableof{\catorder-2}$}

            \draw (-2,-1) rectangle (0,3);
            \drawtextnode{-1}{1}{\corelabelsize $F^{\catorder-1}$}

            \draw (-4,-2) -- (5,-2);
            \drawsmalltensor{6}{-2}{$\hgate$}
            \drawhwire{7}{-2}{$\headvariableof{0}$}
            \draw (3,-2) -- (3,-1);
            \drawvariabledot{3}{-2}

            \drawhwire{0}{2.5}{$\headvariableof{\catorder-1}$}
            \drawtextnode{1}{1.5}{\colorlabelsize $\vdots$}
            \drawhwire{0}{-0.5}{$\headvariableof{1}$}
            \draw (2,-1) rectangle (4,3);
            \drawtextnode{3}{1}{\corelabelsize $U^{\catorder-1}$}
            \draw (4,2.5) -- (7,2.5);
            \drawtextnode{5}{1.5}{\colorlabelsize $\vdots$}
            \drawhwire{7}{2.5}{$\headvariableof{\catorder-1}$}
            \draw (4,-0.5) -- (7,-0.5);
            \drawhwire{7}{-0.5}{$\headvariableof{1}$}

        \end{scope}

    }
    By iteratively decomposing $F^{\catorder-1}$ until we arrive at $F^{1}$ (which is $\hgate$), we get a decomposition of $F^{\catorder}$ into $\catorder$ $\hgate$ matrices and controlled $U$ unitaries.

    \section{Fast Fourier Transform}

    \subsection{Layerwise contraction}

    We contract the FFT tensor network layerwise.
    There are $2\cdot \catorder$ layers, each has a contraction demand in $\mathcal{O}(2^{\catorder})$ (exploiting diagonality of $U$).
    The algorithm has a total computational demand in $\mathcal{O}(d2^{\catorder})$.

    \subsection{Divide and conquer}

    Called the radix-2 FFT \cite{cooley_algorithm_1965}.
    Can be seen as an iterative simplification by adding a $\onesat{\catvariableof{\catorder-1}}$ and splitting it into the basis vectors.


    \stantikzpic{

        \draw (-5,-3) rectangle (-3,3);
        \drawtextnode{-4}{0}{\corelabelsize $\exvector$}

        \drawhwire{-3}{2.5}{$\catvariableof{0}$}
        \drawtextnode{-2}{0.5}{\colorlabelsize $\vdots$}
        \drawhwire{-3}{-2.5}{$\catvariableof{\catorder-1}$}
        \draw (-1,-3) rectangle (1,3);
        \drawtextnode{0}{0}{\corelabelsize $F^{\catorder}$}
        \drawhwire{1}{2.5}{$\headvariableof{\catorder-1}$}
        \drawtextnode{2}{0.5}{\colorlabelsize $\vdots$}
        \drawhwire{1}{-2.5}{$\headvariableof{0}$}

        \drawtextnode{6}{0}{$=$}

        \begin{scope}[shift={(15,0)}]

            \draw (-6,-3) rectangle (-4,3);
            \drawtextnode{-5}{0}{\corelabelsize $\exvector$}

            \drawhwire{-4}{2.5}{$\catvariableof{0}$}
            \drawtextnode{-3}{1.5}{\colorlabelsize $\vdots$}
            \drawhwire{-4}{-0.5}{$\catvariableof{\catorder-2}$}

            \draw (-2,-1) rectangle (0,3);
            \drawtextnode{-1}{1}{\corelabelsize $F^{\catorder-1}$}

            \drawvariabledot{-1}{-2}
            \draw (-1,-2) -- (-1,-3);
            \drawtensorblock{-1}{-4}{2}{$\fbasis+\tbasis$}

            \draw (-4,-2) -- (5,-2);
            \drawsmalltensor{6}{-2}{$\hgate$}
            \drawhwire{7}{-2}{$\headvariableof{0}$}
            \draw (3,-2) -- (3,-1);
            \drawvariabledot{3}{-2}

            \drawhwire{0}{2.5}{$\headvariableof{\catorder-1}$}
            \drawtextnode{1}{1.5}{\colorlabelsize $\vdots$}
            \drawhwire{0}{-0.5}{$\headvariableof{1}$}
            \draw (2,-1) rectangle (4,3);
            \drawtextnode{3}{1}{\corelabelsize $U^{\catorder-1}$}
            \draw (4,2.5) -- (7,2.5);
            \drawtextnode{5}{1.5}{\colorlabelsize $\vdots$}
            \drawhwire{7}{2.5}{$\headvariableof{\catorder-1}$}
            \draw (4,-0.5) -- (7,-0.5);
            \drawhwire{7}{-0.5}{$\headvariableof{1}$}

        \end{scope}

    }

    Now by linearity of contractions, this splits into two terms

    \stantikzpic{

        \begin{scope}[shift={(0,0)}]

            \draw (-6,-3) rectangle (-4,3);
            \drawtextnode{-5}{0}{\corelabelsize $\exvector$}

            \drawhwire{-4}{2.5}{$\catvariableof{0}$}
            \drawtextnode{-3}{1.5}{\colorlabelsize $\vdots$}
            \drawhwire{-4}{-0.5}{$\catvariableof{\catorder-2}$}

            \draw (-2,-1) rectangle (0,3);
            \drawtextnode{-1}{1}{\corelabelsize $F^{\catorder-1}$}

            \draw (-4,-2.5) -- (-3,-2.5);
            \drawsmalltensor{-2}{-2.5}{$\fbasis$}

            \drawsmalltensor{1}{-2.5}{$\fbasis$}

            \draw (2,-2.5) -- (5,-2.5);
            \drawsmalltensor{6}{-2.5}{$\hgate$}
            \drawhwire{7}{-2.5}{$\headvariableof{0}$}
            \draw (3,-2.5) -- (3,-1);
            \drawvariabledot{3}{-2.5}

            \drawhwire{0}{2.5}{$\headvariableof{\catorder-1}$}
            \drawtextnode{1}{1.5}{\colorlabelsize $\vdots$}
            \drawhwire{0}{-0.5}{$\headvariableof{1}$}
            \draw (2,-1) rectangle (4,3);
            \drawtextnode{3}{1}{\corelabelsize $U^{\catorder-1}$}
            \draw (4,2.5) -- (7,2.5);
            \drawtextnode{5}{1.5}{\colorlabelsize $\vdots$}
            \drawhwire{7}{2.5}{$\headvariableof{\catorder-1}$}
            \draw (4,-0.5) -- (7,-0.5);
            \drawhwire{7}{-0.5}{$\headvariableof{1}$}

        \end{scope}

        \drawtextnode{15}{0}{$+$}

        \begin{scope}[shift={(25,0)}]

            \draw (-6,-3) rectangle (-4,3);
            \drawtextnode{-5}{0}{\corelabelsize $\exvector$}

            \drawhwire{-4}{2.5}{$\catvariableof{0}$}
            \drawtextnode{-3}{1.5}{\colorlabelsize $\vdots$}
            \drawhwire{-4}{-0.5}{$\catvariableof{\catorder-2}$}

            \draw (-2,-1) rectangle (0,3);
            \drawtextnode{-1}{1}{\corelabelsize $F^{\catorder-1}$}

            \draw (-4,-2.5) -- (-3,-2.5);
            \drawsmalltensor{-2}{-2.5}{$\tbasis$}

            \drawsmalltensor{1}{-2.5}{$\tbasis$}

            \draw (2,-2.5) -- (5,-2.5);
            \drawsmalltensor{6}{-2.5}{$\hgate$}
            \drawhwire{7}{-2.5}{$\headvariableof{0}$}
            \draw (3,-2.5) -- (3,-1);
            \drawvariabledot{3}{-2.5}

            \drawhwire{0}{2.5}{$\headvariableof{\catorder-1}$}
            \drawtextnode{1}{1.5}{\colorlabelsize $\vdots$}
            \drawhwire{0}{-0.5}{$\headvariableof{1}$}
            \draw (2,-1) rectangle (4,3);
            \drawtextnode{3}{1}{\corelabelsize $U^{\catorder-1}$}
            \draw (4,2.5) -- (7,2.5);
            \drawtextnode{5}{1.5}{\colorlabelsize $\vdots$}
            \drawhwire{7}{2.5}{$\headvariableof{\catorder-1}$}
            \draw (4,-0.5) -- (7,-0.5);
            \drawhwire{7}{-0.5}{$\headvariableof{1}$}

        \end{scope}

    }

    We simplify them to


    \stantikzpic{

        \draw (-6,-3) rectangle (-4,3);
        \drawtextnode{-5}{0}{\corelabelsize $\exvector$}

        \drawhwire{-4}{2.5}{$\catvariableof{0}$}
        \drawtextnode{-3}{1.5}{\colorlabelsize $\vdots$}
        \drawhwire{-4}{-0.5}{$\catvariableof{\catorder-2}$}

        \draw (-2,-1) rectangle (0,3);
        \drawtextnode{-1}{1}{\corelabelsize $F^{\catorder-1}$}

        \draw (-4,-2.5) -- (-3,-2.5);
        \drawsmalltensor{-2}{-2.5}{$\fbasis$}

        \drawsmalltensor{1}{-2.5}{$\hplusbas$}


        \drawhwire{2}{-2.5}{$\headvariableof{0}$}

        \draw (0,2.5) -- (2,2.5);
        \drawhwire{2}{2.5}{$\headvariableof{\catorder-1}$}
        \drawtextnode{1}{1.5}{\colorlabelsize $\vdots$}
        \draw (0,-0.5) -- (2,-0.5);
        \drawhwire{2}{-0.5}{$\headvariableof{1}$}


        \drawtextnode{6}{0}{$+$}

        \begin{scope}[shift={(15,0)}]

            \draw (-6,-3) rectangle (-4,3);
            \drawtextnode{-5}{0}{\corelabelsize $\exvector$}

            \drawhwire{-4}{2.5}{$\catvariableof{0}$}
            \drawtextnode{-3}{1.5}{\colorlabelsize $\vdots$}
            \drawhwire{-4}{-0.5}{$\catvariableof{\catorder-2}$}

            \draw (-2,-1) rectangle (0,3);
            \drawtextnode{-1}{1}{\corelabelsize $F^{\catorder-1}$}

            \draw (-4,-2.5) -- (-3,-2.5);
            \drawsmalltensor{-2}{-2.5}{$\tbasis$}

            \drawsmalltensor{1}{-2.5}{$\hminusbas$}
            \draw (2,-2.5) -- (4,-2.5);
            \drawhwire{4}{-2.5}{$\headvariableof{0}$}


            \drawhwire{0}{2.5}{$\headvariableof{\catorder-1}$}
            \drawtextnode{1}{1.5}{\colorlabelsize $\vdots$}
            \drawhwire{0}{-0.5}{$\headvariableof{1}$}
            \draw (2,-1) rectangle (4,3);
            \drawtextnode{3}{1}{\corelabelsize $U^{\catorder-1}$}

            \drawtextnode{5}{1.5}{\colorlabelsize $\vdots$}
            \drawhwire{4}{2.5}{$\headvariableof{\catorder-1}$}
            \drawhwire{4}{-0.5}{$\headvariableof{1}$}

        \end{scope}

    }

    The contraction of $\exvector$ with $\tbasisat{\catvariableof{\catorder-1}}$ respectively $\tbasisat{\catvariableof{\catorder-1}}$ are the even and the odd coordinates of $\exvector$.
    Both are Fourier transformed with ${\catorder-1}$ bits, which are further divided into Fourier transforms of lower order.
    In the conquer step of the FFT, the odd fourier transform gets a coordinatewise transform by $U$, which is efficient since $U$ is diagonal.


    \section{Quantum Fourier Transform}

    When $\vectorat{\shortcatvariables}$ is a quantum state, the Fourier tansform can be understood as a unitary transforming it.
    Since both the controlled $U$ tensor and the Hadamard $H$ are unitary, each of them can be implemented by a sequence of gates.
    For sparse implementation one can exploit that $U$ is a tensor product of matrixes (i.e. single qubit gates).

    \bibliographystyle{plainnat}
    \bibliography{../../references.bib}


\end{document}